% Section 22.1: 01-riemann-tensor | Part: 4 | Depends on: sec:conn, sec:lc | Status: done | Original lines: 3155-3177
%!TEX root = ../../../main.tex
\begin{intuition}[{curvature = failure of parallel transport to be path independent}]
On a flat plane, if you parallel-transport a vector from $p$ to $q$ along two different paths you obtain the same vector. On a sphere you do not (Example~\ref{ex:spheretransport}: the octant triangle rotates the vector by $\pi/2$). Curvature measures this effect \emph{infinitesimally}: carry $v$ around a tiny parallelogram with sides $\varepsilon X$ and $\delta Y$; it comes back changed by
\[
\Delta v\ \approx\ \varepsilon\,\delta\; R(X,Y)v
\]
(up to a sign depending on the orientation of the loop). Algebraically, curvature is the failure of covariant derivatives to commute: $\nabla_X\nabla_Y-\nabla_Y\nabla_X\neq\nabla_{[X,Y]}$. In GR, ``gravity'' is precisely curvature, the part that cannot be removed by changing coordinates.
\end{intuition}

\begin{formal}[{Riemann curvature}]
For any connection $\nabla$ and fields $X,Y,Z$ define
\[
\boxed{\ R(X,Y)Z=\nabla_X\nabla_YZ-\nabla_Y\nabla_XZ-\nabla_{[X,Y]}Z\ }
\]
This expression is $C^\infty(M)$-linear in each of $X,Y,Z$, so it defines a tensor $R\in\TF{1}{3}(M)$, the \emph{Riemann curvature tensor}. In a chart, $R=R^l{}_{ijk}\,\dd x^i\otimes\dd x^j\otimes\dd x^k\otimes\partial_l$, i.e.\ $R(\partial_i,\partial_j)\partial_k=R^l{}_{ijk}\partial_l$, with
\begin{equation}\label{eq:riemanncomp}
\boxed{\ R^l{}_{ijk}=\partial_i\Gamma^l{}_{jk}-\partial_j\Gamma^l{}_{ik}+\Gamma^s{}_{jk}\Gamma^l{}_{is}-\Gamma^s{}_{ik}\Gamma^l{}_{js}\ }
\end{equation}
(derive it by expanding $\nabla_{\partial_i}\nabla_{\partial_j}\partial_k-\nabla_{\partial_j}\nabla_{\partial_i}\partial_k$). 
\end{formal}
\begin{pitfall}[{index conventions}]
% REWRITTEN: "of the notes" -> "here"; the mention of a specific set of lecture notes is generalised (original wording in legacy/original.tex).
Authors place the indices differently. In Carroll's book \cite{Carroll} $R^\rho{}_{\sigma\mu\nu}$ with $\rho$ the output, $\sigma$ the input vector and $(\mu\nu)$ the pair of directions: it equals $R^l{}_{ijk}$ used here with $(l,k,i,j)=(\rho,\sigma,\mu,\nu)$. The sign of the whole tensor can also differ (some lecture notes use the opposite overall sign). Always check the definition before comparing formulas.
\end{pitfall}
